\documentclass[11pt]{ltjsarticle}
\usepackage{luatexja-fontspec}
\usepackage{amsmath,amssymb}
\usepackage[margin=22mm]{geometry}
\usepackage{graphicx}
\usepackage{longtable,booktabs,array}
\usepackage{calc}
\usepackage{xcolor}
\usepackage{hyperref}
\hypersetup{colorlinks=true,linkcolor=blue,urlcolor=blue,citecolor=blue}
\makeatletter
\def\maxwidth{\ifdim\Gin@nat@width>\linewidth\linewidth\else\Gin@nat@width\fi}
\makeatother
\setkeys{Gin}{width=\maxwidth,keepaspectratio}
\setcounter{secnumdepth}{0}
\setlength{\parindent}{0pt}
\setlength{\parskip}{0.5em}
\providecommand{\tightlist}{\setlength{\itemsep}{0pt}\setlength{\parskip}{0pt}}
\providecommand{\pandocbounded}[1]{#1}
\providecommand{\real}[1]{#1}
\newcounter{none}
\begin{document}
\section{Appendix: On the Emergence of Multidimensional Structure from
the Logarithmic Fundamental
Quantity}\label{appendix-on-the-emergence-of-multidimensional-structure-from-the-logarithmic-fundamental-quantity}

\subsection{--- A Formal Observation on the Dimension of the Hyperplane
Defined by the Additive
Identity}\label{a-formal-observation-on-the-dimension-of-the-hyperplane-defined-by-the-additive-identity}

\textbf{Author:} Noriaki Kihara (ORCID: 0009-0004-6753-4020)
\textbf{Date:} 18 June 2026 \textbf{DOI (this version):}
10.5281/zenodo.20741713 \textbf{Concept DOI (all versions):}
10.5281/zenodo.20741712 \textbf{Zenodo:}
https://zenodo.org/records/20741713 \textbf{Related paper:} ``On the
Logarithmic Map of Scale and the Separation of Scale and Quantization in
a High-Symmetry Limit,'' §3, §5, §7 (Concept DOI:
\href{https://doi.org/10.5281/zenodo.20740841}{10.5281/zenodo.20740841})

\begin{center}\rule{0.5\linewidth}{0.5pt}\end{center}

\subsection{Purpose}\label{purpose}

The related paper stated that structure arises under the strong
constraints of high symmetry and anonymity. In the main text, however,
it did not exhibit that ``arising structure'' in a concrete form. This
appendix is a \textbf{formal observation} intended to fill that gap.

Namely, from the minimal starting point of the fundamental quantity
N=log\textsubscript{2}ν and the additive identity of the conjugate relation alone,
increasing the number of variables makes a multidimensional hyperplane
structure arise formally --- this is shown, as an example, by increasing
the number of variables up to four.

\textbf{What this appendix does not claim (stated up front):}

\begin{itemize}
\tightlist
\item
  That this connects to physical spacetime / dimensions. Not claimed at
  all.
\item
  That four dimensions are special, or that things stabilize at four.
  Not claimed. Four is merely an example and is given no further
  meaning. Asking why a particular dimension would rapidly make the
  discussion difficult, so this appendix does not enter there.
\item
  That this is a mechanism of dimensional emergence. Not claimed.
\end{itemize}

This appendix is confined to an illustration of structure: ``from the
minimal logarithmic structure, a multidimensional hyperplane can be
formally derived according to the number of variables.''

\begin{center}\rule{0.5\linewidth}{0.5pt}\end{center}

\subsection{1. Starting point: the fundamental quantity N (one
variable)}\label{starting-point-the-fundamental-quantity-n-one-variable}

As in §4 of the related paper, the fundamental quantity is introduced as
the logarithm of a squared quantity:

\[N \equiv \tfrac{1}{2}\log_2(\nu^2) = \log_2\nu .\]

When ν=2\textsuperscript{n}, N=n (integer). N appears as a single axis (the whole integer
set $\mathbb{Z}$). This is, on the face of it, a one-dimensional structure.

The question of this appendix is: when this fundamental quantity is
placed in several copies (n of them), does a higher-dimensional
hyperplane structure arise formally according to the number of
variables? The multidimensionality does not ``spring from nothing out of
a single axis''; it comes from the number n of variables placed --- this
is stated up front.

\begin{center}\rule{0.5\linewidth}{0.5pt}\end{center}

\subsection{2. The hyperplane defined by the additive
identity}\label{the-hyperplane-defined-by-the-additive-identity}

As in §5 of the related paper, taking the base-2 logarithm of the
conjugate relation νλ=k turns the product into a sum:

\[\nu\lambda=k \;\xrightarrow{\ \log_2\ }\; N_\nu + N_\lambda = K \quad(K\equiv\log_2 k).\]

We now generalize. Take n fundamental quantities
\(N_1, N_2, \dots, N_n\) and require their sum to be held at a constant
value K:

\[\sum_{i=1}^{n} N_i = K .\]

This is a single linear equation in n variables, and it defines an
\textbf{(n−1)-dimensional hyperplane} within the n-dimensional
coordinate space \((N_1,\dots,N_n)\).

That is, the additive identity raises a flat structure (a hyperplane) of
dimension (n−1) according to the number of variables n.~Below we look
concretely at n=2, 3, 4.

\textbf{Positioning within the axiom system (important):} Placing
several variables (n≥3) is itself an operation beyond Axiom 2 of the
related paper (``only up to 2 can be distinguished in the interior''),
and \textbf{presupposes the dimensional generation of §3} (the genesis
of an integer mode number n via boundedness = the breaking of
anonymity). That is, this appendix shows that, in the regime after
anonymity has been broken and dimension has arisen, the additive
identity defines an (n−1)-dimensional hyperplane. §3 (the genesis of
dimension) and this appendix (n variables → an (n−1)-dimensional
hyperplane) stand in an upstream/downstream relation.

\begin{center}\rule{0.5\linewidth}{0.5pt}\end{center}

\subsection{3. n=2: a line
(one-dimensional)}\label{n2-a-line-one-dimensional}

The sum of two variables \(N_1, N_2\) is constant:

\[N_1 + N_2 = K .\]

This is a \textbf{line} of slope −1 in the \((N_1, N_2)\) plane
(one-dimensional). The relation ``the sum of the conjugate pair is
constant'' appears as a single line in the plane (Figure 1).

\begin{figure}
\centering
\pandocbounded{\includesvg[keepaspectratio]{figure1_n2_line.pdf}}
\caption{Figure 1: n=2. The line of slope −1 defined by the additive
identity N\textsubscript{1}+N\textsubscript{2}=K (1-dimensional)}
\end{figure}

\begin{center}\rule{0.5\linewidth}{0.5pt}\end{center}

\subsection{4. n=3: a plane
(two-dimensional)}\label{n3-a-plane-two-dimensional}

The sum of three variables \(N_1, N_2, N_3\) is constant:

\[N_1 + N_2 + N_3 = K .\]

This is a \textbf{plane} in the three-dimensional coordinate space
\((N_1, N_2, N_3)\) (two-dimensional). In the positive region of the
three axes, this plane appears as a triangular region (Figure 2).

\begin{figure}
\centering
\pandocbounded{\includesvg[keepaspectratio]{figure2_n3_plane.pdf}}
\caption{Figure 2: n=3. The plane defined by the additive identity
N\textsubscript{1}+N\textsubscript{2}+N\textsubscript{3}=K (2-dimensional). The positive region is a triangle
(2-simplex)}
\end{figure}

\begin{center}\rule{0.5\linewidth}{0.5pt}\end{center}

\subsection{5. n=4: a 3-dimensional hyperplane
(three-dimensional)}\label{n4-a-3-dimensional-hyperplane-three-dimensional}

The sum of four variables \(N_1, N_2, N_3, N_4\) is constant:

\[N_1 + N_2 + N_3 + N_4 = K .\]

This is a \textbf{3-dimensional hyperplane} in four-dimensional
coordinate space. Since four dimensions cannot be drawn directly, it is
shown as a projection into three dimensions. In the positive region of
the four axes, this hyperplane appears as a tetrahedral (simplex) region
(Figures 3 and 4).

\begin{figure}
\centering
\pandocbounded{\includesvg[keepaspectratio]{figure3_n4_hyperplane.pdf}}
\caption{Figure 3: n=4. The 3-dimensional hyperplane defined by the
additive identity N\textsubscript{1}+N\textsubscript{2}+N\textsubscript{3}+N\textsubscript{4}=K (3-dimensional). The positive region is
a tetrahedron (3-simplex); the vertices are the four axis intercepts}
\end{figure}

Increasing the number of variables to four here is merely an example. No
special meaning is given to the number four. Increasing the number of
variables further raises an (n−1)-dimensional hyperplane in the same
manner.

\begin{center}\rule{0.5\linewidth}{0.5pt}\end{center}

\subsection{6. Summary}\label{summary}

From the minimal logarithmic fundamental quantity N=log\textsubscript{2}ν and the
additive identity ΣN\textsubscript{i}=K alone, the following structure arises formally
according to the number of variables n:

{\def\LTcaptype{none} % do not increment counter
\begin{longtable}[]{@{}
  >{\raggedright\arraybackslash}p{(\linewidth - 6\tabcolsep) * \real{0.2500}}
  >{\raggedright\arraybackslash}p{(\linewidth - 6\tabcolsep) * \real{0.2500}}
  >{\raggedright\arraybackslash}p{(\linewidth - 6\tabcolsep) * \real{0.2500}}
  >{\raggedright\arraybackslash}p{(\linewidth - 6\tabcolsep) * \real{0.2500}}@{}}
\toprule\noalign{}
\begin{minipage}[b]{\linewidth}\raggedright
Number of variables n
\end{minipage} & \begin{minipage}[b]{\linewidth}\raggedright
Additive identity
\end{minipage} & \begin{minipage}[b]{\linewidth}\raggedright
Arising structure
\end{minipage} & \begin{minipage}[b]{\linewidth}\raggedright
Dimension
\end{minipage} \\
\midrule\noalign{}
\endhead
\bottomrule\noalign{}
\endlastfoot
1 & (a single axis) & axis & 1 \\
2 & N\textsubscript{1}+N\textsubscript{2}=K & line & 1 \\
3 & N\textsubscript{1}+N\textsubscript{2}+N\textsubscript{3}=K & plane & 2 \\
4 & N\textsubscript{1}+N\textsubscript{2}+N\textsubscript{3}+N\textsubscript{4}=K & 3-D hyperplane & 3 \\
n & ΣN\textsubscript{i}=K & hyperplane & n−1 \\
\end{longtable}
}

\begin{figure}
\centering
\pandocbounded{\includesvg[keepaspectratio]{figure4_summary.pdf}}
\caption{Figure 4: n=1 (axis) → n=2 (line) → n=3 (plane) → n=4
(tetrahedron). The dimension of the hyperplane defined by the additive
identity ΣN\textsubscript{i}=K rises as (n−1) with the number of variables n}
\end{figure}

Placing the fundamental quantity N=log\textsubscript{2}ν in several copies (n of them)
and, from the additive identity ΣN\textsubscript{i}=K that holds their sum constant
alone, an (n−1)-dimensional hyperplane arises formally according to the
number n of variables. The multidimensionality does not spring from a
single axis; it comes from the number of variables placed. This is the
minimal concrete example of what the related paper meant by ``structure
arises.''

These hyperplanes are all flat (carry no curvature). This is consistent
with the observation in §7 of the related paper that ``the logarithmic
map sends the sum of squares to a flat linear sum (a hyperplane).''

\begin{center}\rule{0.5\linewidth}{0.5pt}\end{center}

\subsection{7. Limitations}\label{limitations}

\begin{enumerate}
\def\labelenumi{\arabic{enumi}.}
\item
  \textbf{Four is merely an example.} Increasing the number of variables
  to four has no special meaning. No claim is made that four dimensions
  stabilize or are special. The appendix does not enter the question of
  why a particular dimension (entering it would rapidly make the
  discussion difficult).
\item
  \textbf{No physical connection is claimed.} It is not claimed that the
  arising hyperplane connects to physical spacetime / dimensions /
  geometry. This appendix is a purely formal illustration of structure
  and does not address isomorphism with a physical theory.
\item
  \textbf{Contraction by full symmetry is not treated.} That under full
  symmetry (all N\textsubscript{i} equal) the hyperplane contracts to a single point (on
  the diagonal) requires a separate discussion of its own and is not
  treated here (it is left to a separate appendix). This appendix shows
  only that the hyperplane arises.
\item
  \textbf{This is an illustration of possibility, not a proof of
  necessity.} It only shows by one example that ``a multidimensional
  structure can arise,'' and does not show that it ``must arise.''
\end{enumerate}

\begin{center}\rule{0.5\linewidth}{0.5pt}\end{center}

\emph{This appendix is a formal illustration of the related paper's
claim (that structure arises under strong constraints), and contains no
definitive claim regarding physical connection, the necessity of a
particular dimension, or stability.}
\end{document}
